\documentclass{article}
\usepackage{tikz}
\usepackage{pgfplots}
\usepackage{amsmath, amsfonts}
\usepackage{enumerate}
\usepackage{makecell}
\usepackage{xypic}
\usepackage{fancyvrb}
\usepackage{listings}
\input{title.tex}

\newcommand{\N}{\mathbb{N}}
\newcommand{\R}{\mathbb{R}}
\begin{document}
    \maketitle
    \section{}
    Wir zeigen zunächst, dass die Harmonische
    Reihe $H_n$ in $O(\log n)$ liegt.
    Dafür definieren wir die Ableitung der Harmonischen
    Reihe. Da Jeder punkt der Harmonischen Reihe
    2 Steigungen zum vorherigen / nächsten Punkt hat,
    wählen wir das Maximum der Beiden $(H_n)':=\frac{1}{n}$.
    Es ist also
    \begin{align*}
        0\leq\lim\limits_{n\rightarrow\infty}
        \frac{H_n}{\log n}&\leq
        \lim\limits_{n\rightarrow\infty}
        \frac{\frac{1}{n}}{\frac{1}{n}}=1\\
    \end{align*}
    Demnach ist $H_n\in O(\log n)$

    \begin{tabular}{cc}
        Zug&Wkeit\\
        1&1\\
        1&$\frac{1}{2}$\\
        $\hdots$\\
        n&$\frac{1}{n}$
    \end{tabular}

    Demnach ist $E=H_n\in O(\log n)$
\end{document}
